%% 
%% Template article for cas-sc documentclass for 
%% double column output.

\documentclass[a4paper,fleqn]{cas-sc}

% If the frontmatter runs over more than one page
% Use the longmktitle option.

%\documentclass[a4paper,fleqn,longmktitle]{cas-sc}

%\usepackage[numbers]{natbib}
%\usepackage[authoryear]{natbib}
\usepackage[numbers]{natbib}
%% The amssymb package provides various useful mathematical symbols
\usepackage{amssymb}
%% The amsmath package provides various useful equation environments.
\usepackage{amsmath}
\usepackage{epstopdf}
\usepackage{mathptmx,mathtools}
\usepackage{subcaption}
\usepackage{float} % Add this in your preamble
\usepackage[export]{adjustbox}
\usepackage{array} % Needed for m{} column type
\usepackage[percent]{overpic} % <-- put this in your preamble



\usepackage{amsmath}
\usepackage{nicefrac}
\usepackage{dirtytalk}
\usepackage{graphicx}   % For including images
\usepackage{tikz}
\usetikzlibrary{calc,tikzmark} % needed <<<<<<<<<<
\usepackage{xcolor}     % For coloring
\usepackage{lineno}
\renewcommand\linenumberfont{\normalfont\bfseries\small} % <-- Increase line number size
\usetikzlibrary{arrows.meta}
\usetikzlibrary{calc, positioning}
%%%Author macros
\def\tsc#1{\csdef{#1}{\textsc{\lowercase{#1}}\xspace}}
\tsc{WGM}
\tsc{QE}

\begin{document}
	
	\setcounter{figure}{1}

			%%%%%%%%%%%%%%%%%%%%%%%%%%%%%________Figure_6___________%%%%%%%%%%%%%%%%%%%%%
	\begin{figure}[!t]
		\centering
		
		% -------- Subfigure (a) -------- %
		\begin{subfigure}[t]{0.556\textwidth}
			\centering
			\begin{overpic}[width=\textwidth,trim={1cm 0cm 0cm 0cm},clip]{Fig5a.eps}
				\put(4,80){\fontfamily{ptm}\selectfont\textbf{(a)}}
				
				
			\end{overpic}
		\end{subfigure}%
		% -------- Subfigure (b) -------- %
		\begin{subfigure}[t]{0.45\textwidth}
			\centering
			\begin{overpic}[width=\textwidth,trim={0.5cm 0cm 2cm 0cm},clip]{Fig5b.eps}
				\put(5,97){\fontfamily{ptm}\selectfont\textbf{(b)}}
				
			\end{overpic}
		\end{subfigure}
		
		\caption{Temporal backward (a–c) and forward (e–g) evolution of the wake and its influence on the surrounding flow regions around the rising bubble for case $C_6$. The flow regions are separated by the hyperbolic LCSs obtained from the forward FTLE field with $\Delta\tau = 10$ at $\tau = 3.6$, as shown in (d), where the outermost and innermost structures are indicated by red and blue solid lines, respectively. The corresponding elliptic LCSs, shown in green with a color spectrum increasing toward the vortex core, highlight regions of rotational coherence. Panels (a–c) correspond to $\tau = 0.05$, $0.75$, and $2.0$, respectively, representing the backward evolution, while panels (e–g) correspond to $\tau = 5.4$, $7.0$, and $8.0$, showing the forward evolution. The light-blue regions represent the surrounding fluid, while the yellow and dark-blue regions denote the fluid above the bubble, with the dark-blue area located near the bubble top surface.}
		
		\label{Fig6}
	\end{figure}
	%%%%%%%%%%%%%%%%%%%%%%%%%%%%%________Figure_6_end___________%%%%%%%%%%%%%%%%%%%%%
	
\end{document}